\section{模型假设}
为在不改变题意的前提下统一四问建模口径，作如下假设。

\begin{enumerate}
\item \textbf{时间离散假设。} 一天划分为144个10min 时段，各时段内的负荷、光伏、电价与储能充放电功率按该时段代表值计算，时段长度 \(\Delta t=1/6\) h。
\item \textbf{时间标签口径。} 全文统一把数据标签解释为对应时段的起始时刻，例如0:10表示区间 \([0{:}10,0{:}20)\)。
\item \textbf{光伏优先供负荷。} 任意时段先利用光伏满足本地负荷，只有负荷满足后的富余光伏才可用于储能充电；若仍无法消纳，则记为弃光。
\item \textbf{不考虑向外网售电。} 题目未给出富余电量上网价格，因此模型仅考虑从外部电网购电，不设置反向售电变量。富余光伏用 \(V\) 表示，计划购电中因预测误差未被实际利用的电量用 \(W\) 表示，两者含义严格区分。
\item \textbf{储能充放电效率。} 充电效率与放电效率均按单程90\% 处理，即 \(\eta_c=\eta_d=0.90\)。
\item \textbf{储能物理边界。} 储能运行电量保持在 \([1200,\ 10800]\) kWh，最大充放电功率均为5000kW。
\item \textbf{充放电互斥。} 同一10min 时段内储能不同时充电和放电，通过二元状态变量 \(u\) 强制实现。
\item \textbf{问题一首末电量约束。} 问题一按题目要求设置 \(E_{144}=E_0=6000\) kWh。
\item \textbf{问题二至问题四跨日连续。} 后三问不再强制每日首末储电量相同，而满足 \(E_{d+1,0}^{act}=E_{d,T}^{act}\)，从而保留储能的跨日调节能力。
\item \textbf{紧急购电用途。} 问题二至问题四中，紧急购电仅用于补足当期实际负荷缺口，不用于储能充电。
\item \textbf{计划购电不可事后无成本撤销。} 当计划购电量超过实际负荷和储能可用需求时，剩余部分记为已购未用电量 \(W\)，不得将其直接从计划中消失或售回外网。
\item \textbf{信息因果性。} 任一决策时刻只能使用该时刻之前已经获得的历史真实数据、已经发布的预测和当前真实储电量；不得提前读取未来真实负荷、未来真实光伏、未来真实价格和尚未发布的预报。
\item \textbf{SAA 历史误差保持日内结构。} 负荷、光伏和价格误差均以“完整历史日”为基本抽样单元，不逐10min 时段独立打乱；当同时构造多类随机变量时，尽量从同一历史日期联合抽取，以保留其历史相关结构。
\item \textbf{一月预热。} 问题二至问题四将2025年1月作为已知数据预热期，用于建立预测历史、误差库及储能连续初值；正式统计窗口从2025年2月1日开始。
\item \textbf{问题四真实费用结算。} 问题四中的预测电价和价格情景仅用于计划决策，最终报送费用统一使用附件4已实现的真实电价重新计算。
\end{enumerate}

上述假设中，问题一的首末储电量相同与问题二至问题四的跨日连续属于不同题目条件，不混用。
